% ID: MAT-ALG-INT-011
% Titolo: Espressione con esponente zero
% Disciplina: Matematica
% Area: Algebra
% Argomento: Numeri interi
% Classe: Prima secondaria di secondo grado
% Difficolta: 3
% Tipo: problema_numerico
% Tag: matematica, algebra, numeri_interi
% Fonte: appunti_matematica_anonimizzati
% Provenienza: materiale fornito dal docente; testo anonimizzato; soluzione redatta per Exergo

\begin{esercizio}
Calcola $\left[(3^2\cdot2^2)^5:6^8\right]-\left\{\left[(16^3)^2:(2^5)^4\right]\cdot2^9\right\}:4^5-(2^4\cdot5^2)^0$.
\end{esercizio}

\begin{soluzione}
Il primo termine è $6^{10}/6^8=36$. Il secondo è $2^{24-20+9-10}=8$. L'ultimo è 1, perché la base è diversa da zero. Risultato: $36-8-1=27$.
\end{soluzione}
